% =========================================================================
% secciones_es/sec05_silicon_performance_benchmarking.tex — Sección V, español
% =========================================================================
\section{Rendimiento computacional, exactitud y benchmarking numérico}
\label{sec:benchmarking}

\subsection{Validez de la medición mediante el determinismo del gemelo digital}
\label{sub:determinism}

La Sección~\ref{sub:lifecycle} estableció que la promoción al estado de producción es atómica y bit a bit idéntica a un Golden Master sellado con SHA-512. Las métricas que siguen describen ese estado de software probado y los artefactos de validación publicados. El Algoritmo~\ref{alg:n1_seal} presenta el pipeline de resiliencia y sellado usado para producir el registro informado.

% ---- Algoritmo 2 ----
\begin{algorithm}[!t]
\caption{Barrido de resiliencia $N$--$1$ con bootstrap de bloques de 24 horas y sellado SHA-512}
\label{alg:n1_seal}
\begin{algorithmic}[1]
\Require estado sellado del solver $\Sigma$, catálogo de contingencias $\mathcal{C}$
         ($|\mathcal{C}| = 215$), horizonte $\mathcal{H}$ (87,360 h; 520 ciclos semanales)
\Ensure registro de resiliencia $\mathcal{R}$, sello $\sigma$
\State $\mathcal{R} \gets \emptyset$
\For{$c \in \mathcal{C}$} \Comment{orden determinista, fail-closed}
   \For{$b \in \text{bootstrap}(\mathcal{H}, 24\,\mathrm{h})$} \Comment{bloques de día completo}
      \State inyectar la contingencia $c$ sobre el bloque $b$
      \State ejecutar el Algoritmo 1 en cada hora del bloque
      \State registrar $\rho$, $\kappa_2$, márgenes de tensión y de carga de rama en $\mathcal{R}[c, b]$
   \EndFor
\EndFor
\State agregar los márgenes sobre la envolvente de bootstrap por contingencia
\State $\sigma \gets \mathrm{SHA512}\text{-chain}(\Sigma, \mathcal{C}, \mathcal{R})$ \Comment{manifiesto desacoplado}
\State \Return $\mathcal{R}$, $\sigma$
\end{algorithmic}
\end{algorithm}

\subsection{Velocidad de convergencia y memoria}
\label{sub:speed}

Sobre hardware comercial estándar (un portátil de 8 hilos de la clase Dell Inspiron 5570), el solver particionado converge cada estado horario de AC de la red nacional completa de 142 barras en \textbf{2.65~ms}, con una huella residente de \textbf{473~MiB} para el horizonte de validación. Estas cifras son mediciones experimentales de referencia obtenidas en la estación de trabajo Linux de 8 hilos del autor mediante su motor de investigación C++/Sparse; sirven como línea base de comparación para futuros solvers de la comunidad, mientras que el arnés abierto liberado reproduce determinísticamente los 41 controles matemáticos y estructurales de la liberación (topología, matrices Ybus, la condensación de Schur a $\delta = 2.788 \times 10^{-14}$~pu y los balances estáticos), no la cifra de tiempo de muro. El horizonte comprende 87,360 horas operativas consecutivas correspondientes a 520 semanas completas de 168 horas (3,640 días de lunes a domingo) dentro del período 2016--2025, preservando la ciclicidad semanal del despacho; los 13 días civiles excluidos de la década (6 días huérfanos de borde más los 7 días de la semana festiva) se omiten por construcción (Sección~\ref{sub:dataprotocol}). El barrido completo se termina en minutos en este equipo, en consistencia con las asíntotas $O(N^3)$-a-$O(\sum_k N_k^3 + N_B^3)$ de la Proposición~\ref{prop:complexity} y con la práctica dispersa de reducción de relleno \cite{george1973}, \cite{saad2003}. La resolución de 2.65~ms mide cribado estático de estados AC; no modela el despeje dinámico de faltas, la respuesta de protecciones ni el control en lazo cerrado. El rendimiento por estado permite barridos estáticos multiescenario de gran volumen bajo las condiciones de modelo y hardware indicadas.

Estas mediciones caracterizan el costo computacional de las comprobaciones de factibilidad AC en régimen permanente. Permiten repetir barridos de parámetros, bootstrap y contingencias en el hardware descrito; no demuestran el despliegue de un sistema dinámico de protección o de control en lazo cerrado.

\subsection{Exactitud y condicionamiento numérico}
\label{sub:accuracy}

Tres magnitudes medidas acotan la fidelidad numérica de la solución condensada. Primero, el residuo fasorial de acople de frontera se evalúa bajo dos métricas complementarias y rigurosas:
(i) el \emph{residuo de reconstrucción en el espacio de estados nodales} $\|\mathbf{v}_{\mathrm{rec}} - \mathbf{v}_{\mathrm{full}}\|_{\infty} = 2.788 \times 10^{-14}\ \mathrm{pu}$ (recomputado desde los CSVs de entrada como $\delta_{\mathrm{recomp}} = 2.70 \times 10^{-14}\ \mathrm{pu}$), el cual mide directamente el error fasorial máximo entre el vector de tensiones nodales recuperado por Schur y la solución monolítica independiente de sistema completo $\mathbf{Y}\mathbf{v}_{\mathrm{full}} = \mathbf{b}$; mientras que corridas exploratorias sobre matrices previas a la consolidación registraron un umbral numérico de $6.35 \times 10^{-12}\ \mathrm{pu}$, la calibración topológica y la consolidación sum-invariante de clones de demanda en subestaciones multivel redujeron la propagación del error numérico en la retro-sustitución directa hasta el piso de precisión de máquina de $10^{-14}\ \mathrm{pu}$; y
(ii) el \emph{residuo algebraico matricial del complemento de Schur} $\|\mathbf{S} - (\mathbf{D} - \mathbf{C}\mathbf{A}^{-1}\mathbf{B})\|_{\infty} < 10^{-16}\ \mathrm{pu}$ ($0.00\ \mathrm{pu}$), el cual confirma el cierre algebraico exacto del operador de frontera previo a la retro-sustitución.
El residuo de estado de $2.788 \times 10^{-14}\ \mathrm{pu}$ se sitúa más de un orden de magnitud por debajo del techo global de redondeo escalado por el número de condición de frontera ($\varepsilon_{\mathrm{mach}}\kappa_2(\mathbf{S}) = 2.22 \times 10^{-16} \times 2{,}489.64 \approx 5.53 \times 10^{-13}$~pu) --- el pequeño residuo no nulo, firma de dos trayectorias de factorización distintas, es un beneficio directo de la estructura ortogonal y bien condicionada de la partición de frontera $\Gamma = 30$ --- y casi seis órdenes de magnitud por debajo de la tolerancia de parada de $10^{-8}$~pu de la cascada. Segunda, el controlador sigmoidal de límites reactivos de la Sección~\ref{sub:regularization} contribuye un error de regularización de $5.32 \times 10^{-12}$ mientras elimina el ciclo límite PV--PQ. Tercera, el número de condición espectral de cada sistema de frontera reducido se monitorea y archiva: $\kappa_2(\mathbf{S}) = 2{,}489.64$ (peor caso medido en la frontera del SENI-142) es el número de condición espectral medido del complemento de Schur de frontera, evidenciando un acoplamiento moderado que se mantiene alejado de la singularidad numérica y permite inversión estable en doble precisión; la maquinaria de umbrales de gobernanza se validó sobre referencias sintéticas antes del despliegue nacional (la Jacobiana de $22\times22$ del IEEE 14 barras se clasifica MARGINAL en $\kappa_2 = 117.7823$), de modo que una degradación del condicionamiento en la red real se detectaría contra una escala calibrada y no contra una ad-hoc.

\subsection{Resiliencia a contingencias sobre la topología de referencia}
\label{sub:contingencies}

La batería de resiliencia evalúa el catálogo $N-1$ de 215 contingencias --- la lista propia de salidas del operador de transmisión, cribada bajo la doctrina de selección automática de ramas de \cite{ejebe1979} --- sobre la topología de referencia, concentrada en el anillo metropolitano de 138~kV cuyas distancias de ruta calibradas (70.5~km, 30.86~km, 22.1~km) fijan la geometría eléctrica, como muestra la Figura~\ref{fig:topologia_seni}. El disparo del complejo de Punta Catalina ($2 \times 376 = 752$~MW), la mayor contingencia común del sistema, se resuelve bajo el criterio preventivo de la Sección~\ref{sub:n1}: la pérdida de potencia se cubre mediante el redespacho de la reserva rodante co-optimizada ($\sum r_g^+ \ge 752$~MW), preservando la factibilidad de tensión y flujo de potencia en estado estacionario post-evento sin deslastre de carga ($P_{\mathrm{shed}} = 0$). El registro de anomalías de la campaña también documenta depresión severa de tensión con agotamiento de margen reactivo --- $\lambda_{\min}(\mathbf{J}_R) < 10^{-2}$ sobre el bloque reactivo reducido \eqref{eq:jr} --- en el lazo JSAURI--Palamara--Haina de 138~kV bajo estados $N-1$ específicos; estos hallazgos se añadieron al catálogo de contingencias y activan la bisección de soporte de shunts de la Sección~\ref{sub:n1}. Todos los márgenes se reportan bajo la envolvente de bootstrap de bloques de 24 horas del Algoritmo~\ref{alg:n1_seal}.

La campaña multitemporal de 87,360 estados y el cribado de 215 contingencias constituyen resultados del motor patent-pending C++/Sparse del autor; el paquete público certifica el modelo de red, los invariantes de resolución condensada y los balances de las ocho contingencias canónicas.

El pase de admisión de cuatro cláusulas es un resultado numérico reportable, no una afirmación narrativa. La Tabla~\ref{tab:worstcase} informa, para los ocho elementos más severos del catálogo --- ordenados por margen de tensión post-evento y por carga térmica ---, la fila completa de despacho y verificación: flujo o generación pre-evento, peor tensión post-evento, peor carga de rama post-evento, desviación de la barra de referencia tras el despliegue de reserva, reserva co-optimizada redespachada y carga no servida. Las filas de Punta Catalina merecen énfasis. El disparo unitario (GEN-PC1, 376~MW) es la peor $N-1$ unitaria; el disparo a nivel de complejo (GEN-PC-COMPLEX, $2 \times 376 = 752$~MW) es el mayor evento de modo común de la isla, cubierto preventivamente por el dimensionamiento de reserva de \eqref{eq:reserve}; ambas filas cierran con toda barra retenida dentro de $[0.95, 1.05]$~pu, peor carga de rama por debajo del 92.4\% de la capacidad térmica, desviación de barra de referencia acotada y carga no servida nula. 

LINE-138-CUM-SPMI es el corredor cuya falla en servicio de 2025 está documentada en el registro hemerográfico nacional \cite{listin2025}; su fila del catálogo demuestra que la red de transmisión modelada soporta la pérdida de ese circuito sin deslastre ($P_{\mathrm{shed}} = 0$) bajo factibilidad AC estática. Este desacople analítico entre capacidad de transporte estática y colapsos nacionales históricos está corroborado por registros periciales primarios: investigaciones técnicas independientes del colapso del 11 de noviembre de 2025 \cite{uasdInformeSPM1_2025} establecieron que el apagón fue iniciado por maniobras de mantenimiento en S/E SPM1, la cual carecía de protección diferencial de barras, disparando una cascada descoordinada de protecciones de distancia que desconectó 10 circuitos de transmisión en 680~ms; el informe pericial ejecutivo del Instituto de Energía de la Universidad Autónoma de Santo Domingo (IEUASD) \cite{uasdInformeSPM1_2025} y del Operador del Sistema (registro horario de telemetría de OC-SENI, ventana 13:15--13:30 \cite{ocseniScadaP15_2025}) registran una frecuencia media de la ventana de 15 minutos (P15) de 55.11~Hz durante el colapso, con el 99.89\% de la ventana fuera de la banda operativa normal. De manera análoga, el apagón general del 16 de mayo de 2015 (1,782.65~MW desconectados) se originó en una falla en barraje de 138~kV en Itabo Gas agravada por fallo de interruptor con tiempo de despeje superior a $10\times$ el nominal \cite{listinItaboGas_2015}. El benchmark empírico confirma que el patrimonio de transmisión dominicano posee suficiente margen térmico y de tensión post-evento bajo $N-1$; las vulnerabilidades históricas corresponden a fallas transitorias, dinámicas y de coordinación de protecciones que pertenecen estrictamente al dominio dinámico.

Las dos filas más escrutadas admiten una lectura física explícita. Para el corredor doble circuito PALA--SCN (conductor ACSR 795~kcmil Drake, hilado 26/7, 988~A continuos a $75^{\circ}$C, $S^{\max} = 236.2$~MVA por circuito y 472.4~MVA para el corredor agregado), el estado pre-evento porta 412.7~MW --- 206.4~MW por circuito, una carga inicial del 87.4\%; tras el disparo de un circuito, el gemelo sobreviviente asume 206.7~MW, el 87.6\% reportado, mientras el excedente se redistribuye por la malla metropolitana: el enlace JSAURI--SCN de 230~kV y 30.86~km de línea ($S_{\max} = 418.3$~MVA, portando 356.2~MW, una carga pre-evento del 85.2\%) absorbe unos 142~MW de flujo desviado --- cerca del 23.7\% de la capacidad nominal de 600~MVA del banco de transformación 345/138~kV XFMR-345-138-JS que asiste la transferencia desde la troncal de 345~kV hacia el anillo de 138~kV --- y el circuito JSAURI--PALA de 138~kV y 70.5~km absorbe unos 64~MW (89.2\%), sin violación térmica en régimen permanente. Para el disparo a nivel de complejo de Punta Catalina, el redespacho de reserva co-optimizada reemplaza exactamente la inyección perdida ($\Delta P_g = 752.0$~MW, desplegada en unidades térmicas, hidroeléctricas y renovables geográficamente remotas), de modo que la desviación residual en la barra de referencia es físicamente el incremento de pérdidas de transporte por reenrutar 752~MW a través del patrimonio de 142 barras: $\Delta P_{\mathrm{slack}} = \Delta P_{\mathrm{loss}} = +24.8$~MW, un incremento del 0.72\% sobre la demanda nacional de referencia de 3,450~MW. En la fila LINE-345-PC-JS, el evento se audita como una contingencia compuesta crítica de transporte y evacuación (salida $N-1$ de línea acoplada al aislamiento simultáneo del complejo térmico de $2 \times 376 = 752$~MW por pérdida de su única vía de evacuación troncal), donde el déficit de 752~MW cierra exactamente mediante dos variables de despacho separadas de la formulación multivector: $\Delta P_g = 687.40$~MW de redespacho térmico/hidráulico convencional (Quisqueya~I--II 180.0, AES Andr\'es CC 145.0, Estrella del Mar~III 110.0, Barahona carbón/CEPP 95.0, hidráulicas Valdesia--Tavera 82.4, reserva solar Monte Plata 75.0) y $\Delta P_{\mathrm{storage}} = 58.50$~MW de descarga rápida de almacenamiento (bombeo reversible Sabaneta/Dos Bocas 35.0~MW; BESS metropolitano 23.5~MW), con el residuo $\Delta P_{\mathrm{slack}} = +6.10$~MW de pérdidas marginales de transporte cerrando el balance en exactamente 752.00~MW (error nulo). El ratio de pérdidas resultante del 0.89\% --- frente al 2.5--3.3\% de los redespachos puramente remotos --- se explica por la inyección del almacenamiento en proximidad eléctrica directa del anillo metropolitano de carga, evitando el tránsito por las troncales de larga distancia de 345/230~kV. Debido a que el corredor LINE-345-PC-JS es la única vía troncal de evacuación masiva del complejo, su disparo aísla la inyección completa hacia el SENI, configurando el evento canónico compuesto de dimensionamiento. La convergencia Newton--Raphson se verificó a $\|\Delta P, \Delta Q\|_{\infty} \le 10^{-8}$~pu en a lo sumo cuatro iteraciones por estado, con balance nodal exacto, dentro de la salida de validación sellada de la suite canónica.

% ---- Tabla VI: contingencias críticas peor caso (blindaje R3, EJE 4) ----
\begin{table*}[!t]
\caption{Contingencias Críticas Peor Caso: Informe de Despacho y Verificación Estática Post-Evento. Selección de los ocho elementos más severos del catálogo por margen de tensión post-evento y carga térmica; el registro completo de 215 elementos se archiva en la salida de validación sellada de la suite canónica de 142 barras. Los flujos pre-evento corresponden al estado de punta de referencia; tensiones y cargas son valores post-evento peor caso sobre barras retenidas y ramas sobrevivientes.}
\label{tab:worstcase}
\centering
\footnotesize
\setlength{\tabcolsep}{4.5pt}
\resizebox{\textwidth}{!}{%
\begin{tabular}{@{}lllrrrrrr@{}}
\toprule
\textbf{ID Contingencia} & \textbf{Activo} & \textbf{Tipo} &
\textbf{Flujo/gen} & \textbf{min $V_i$} & \textbf{máx $|S|/S_{\max}$} &
\textbf{$\Delta P_{\mathrm{slack}}$} & \textbf{$\Delta P_g$} & \textbf{$P_{\mathrm{shed}}$} \\
 & & & \textbf{pre-evento (MW)} & \textbf{(pu)} & \textbf{(\%)} &
\textbf{(MW)} & \textbf{(MW)} & \textbf{(MW)} \\
\midrule
GEN-PC1 & Punta Catalina Unidad 1 (376 MW) & $N-1$ unitaria & 376.0 & 0.9612 & 91.7 & $+12.4$ & 376.0 & 0.00 \\
GEN-PC-COMPLEX & Complejo Punta Catalina ($2 \times 376$ MW) & Modo común & 752.0 & 0.9581 & 92.4 & $+24.8$ & 752.0 & 0.00 \\
LINE-345-PC-JS & Troncal 345 kV Punta Catalina--JSAURI & Compuesto (Línea + Aislamiento 752 MW) & 748.3 & 0.9703 & 89.2 & $+6.1$ & 687.4 & 0.00 \\
XFMR-345-138-JS & Banco autotransformador JSAURI 345/138 kV (600 MVA) & $N-1$ de transformador & 521.7 & 0.9748 & 88.5 & $+4.7$ & 186.2 & 0.00 \\
LINE-138-CUM-SPMI & 138 kV SPM--Cumayasa (falla en servicio 2025 \cite{listin2025}) & $N-1$ de línea & 172.1 & 0.9802 & 84.1 & $+2.9$ & 96.4 & 0.00 \\
CORR-138-PALA-SCN & Corredor doble circuito 138 kV PALA--SCN, 22.1 km (pérdida de un circuito) & $N-1$ de corredor & 206.4 & 0.9755 & 87.6 & $+5.2$ & 158.3 & 0.00 \\
LINE-230-JS-SCN & Enlace 230 kV JSAURI--SCN & $N-1$ de línea & 356.2 & 0.9790 & 83.3 & $+3.8$ & 141.7 & 0.00 \\
GEN-AES-ANDRES & Ciclo combinado AES Andr\'es (319 MW) & $N-1$ unitaria & 319.0 & 0.9846 & 78.5 & $+8.3$ & 319.0 & 0.00 \\
\midrule
\multicolumn{9}{@{}p{0.98\textwidth}@{}}{\scriptsize Todas las filas satisfacen el pase de admisión de cuatro cláusulas de la Sección~\ref{sub:n1}: convergencia $\|\Delta P, \Delta Q\|_{\infty} \le 10^{-8}$~pu; $V_i \in [0.95, 1.05]$~pu; $|S_\ell| \le S_\ell^{\max}$; $|\Delta P_{\mathrm{slack}}| \le \varepsilon_{\mathrm{slack}}$ con $P_{\mathrm{shed}} = 0$. $\Delta P_g$ es la reserva co-optimizada desplegada tras el evento; $\Delta P_{\mathrm{slack}}$ es la desviación residual de la barra de referencia tras ese despliegue. Estados de verificación estática de factibilidad AC post-contingencia computados sobre el modelo canónico completo de 142 barras de transmisión bajo el despacho multitemporal base. El disparo GEN-PC-COMPLEX ($2 \times 376 = 752$~MW) representa el caso de referencia para el dimensionamiento de reserva rodante ($R^{\mathrm{req}}$), externo al catálogo de 215 salidas simples ($N-1$). La fila LINE-345-PC-JS se audita como una contingencia compuesta crítica de transporte y evacuación ($N-1$ de línea acoplada al aislamiento simultáneo del complejo térmico de $2 \times 376 = 752$~MW por pérdida de su única troncal de evacuación), distinguiéndose de las salidas puras de líneas de transmisión malladas que redistribuyen flujo sin pérdida de generación (las cuales exhiben $|\Delta P_{\mathrm{slack}}| \le +6.10$~MW a lo largo de las 135 líneas). Su carga pre-evento se descompone como 748.30~MW de flujo troncal más 3.70~MW de servicios auxiliares locales, es decir, 752.00~MW de inyección bruta aislada; los ratios duales canónicos del residuo de $+6.10$~MW en barra de referencia son 0.887\% sobre el redespacho de 687.40~MW y 0.811\% sobre el déficit bruto de 752.00~MW. Para el corredor doble circuito PALA--SCN, el flujo pre-evento es el portado por el circuito disparado (206.4~MW) dentro de un flujo agregado del corredor de 412.7~MW; la carga post-evento reportada (87.6\%) corresponde al circuito sobreviviente y a los enlaces adyacentes, dentro de las capacidades térmicas reales de conductor. Conductor: ACSR 795~kcmil Drake (26/7), 988~A continuos a $75^{\circ}$C, $S^{\max} = 236.2$~MVA por circuito (472.4~MVA agregados).}
\end{tabular}}%
\end{table*}

% ---- Tabla V: Resumen estadístico agregado de las 215 contingencias N-1 tamizadas ----
\begin{table*}[!t]
\caption{Resumen Estadístico Agregado de la Batería Tamizada de 215 Contingencias $N-1$ y Referencia de Modo Común}
\label{tab:contingency_aggregate}
\centering
\footnotesize
\begin{tabular}{@{}lrrrrrr@{}}
\toprule
\textbf{Categoría de Contingencia} & \textbf{Cantidad} & \textbf{min $V_i$ Post-Evento} & \textbf{Peor Carga de Rama} & \textbf{Máx $|\Delta P_{\mathrm{slack}}|$} & \textbf{Total $P_{\mathrm{shed}}$} & \textbf{Violaciones} \\
 & & \textbf{[pu]} & \textbf{[\%]} & \textbf{[MW]} & \textbf{[MW]} & \textbf{(PASS / Total)} \\
\midrule
Líneas de transmisión (345, 230, 138 kV) & 135 & $[0.9521,\, 0.9892]$ & $[32.4\%,\, 89.2\%]$ & $+6.10$ & $0.00$ & $0 / 135\ \text{(PASS)}$ \\
Transformadores de potencia de subestación & 58 & $[0.9550,\, 0.9910]$ & $[28.1\%,\, 88.7\%]$ & $+5.10$ & $0.00$ & $0 / 58\ \text{(PASS)}$ \\
Unidades generadoras despachables principales & 22 & $[0.9612,\, 0.9850]$ & $[45.0\%,\, 91.7\%]$ & $+12.40$ & $0.00$ & $0 / 22\ \text{(PASS)}$ \\
\midrule
\textbf{Total Catálogo Tamizado } $N-1$ & \textbf{215} & $\mathbf{0.9521}$ & $\mathbf{91.7\%}$ & $\mathbf{+12.40}$ & $\mathbf{0.00}$ & $\mathbf{0 / 215\ (PASS)}$ \\
\midrule
Referencia benchmark de modo común ($N-2$) & 1 & $0.9581$ & $92.4\%$ & $+24.80$ & $0.00$ & $0 / 1\ \text{(PASS)}$ \\
\bottomrule
\end{tabular}
\end{table*}

La Tabla~\ref{tab:contingency_aggregate} reporta la distribución estadística agregada a lo largo de la batería completa de 215 contingencias unitarias tamizadas $N-1$ y el caso de referencia externo de modo común. Dentro del catálogo $N-1$ unitario, la desviación máxima de potencia activa en la barra de referencia es $\Delta P_{\mathrm{slack}} = +12.40$~MW (bajo el disparo unitario GEN-PC1), exhibiendo un holgado margen de seguridad de $12.60$~MW por debajo del límite operativo $\varepsilon_{\mathrm{slack}} = 25.0$~MW (un margen de seguridad del $50.4\%$). La desviación de $+24.80$~MW se alcanza exclusivamente bajo el evento de referencia de modo común $N-2$ (GEN-PC-COMPLEX, 752~MW), en el cual la totalidad del déficit de 752~MW se cubre mediante redespacho remoto y la barra slack suministra únicamente el incremento no lineal de pérdidas de transmisión ($3.3\%$). Crucialmente, la ausencia de deslastre de carga ($P_{\mathrm{shed}} = 0$) en las 215 contingencias es un resultado empírico y computado de la resolución primal del flujo AC y del AC-OPF: al hallarse explícitamente un punto de operación admisible sin deslastre para cada contingencia, la penalización de $c^{\mathrm{shed}} = 10{,}000\ \$/\mathrm{MWh}$ garantiza que el optimizador nunca seleccione deslastre de demanda.

% ---- Figura 6 ----
\begin{figure*}[!t]
\centering
\resizebox{0.85\textwidth}{!}{% Figure 6 tikzpicture — Edición en español (misma geometría
% anti-colisión con rótulos de arista con relleno blanco; paleta seni-* del
% preámbulo hospedador).
\begin{tikzpicture}[
    subestacion/.style={circle, draw=seni-blue, fill=seni-blue!10, thick, minimum size=13mm, font=\bfseries\footnotesize, align=center},
    sub345/.style={circle, draw=seni-red, fill=seni-red!10, thick, minimum size=13mm, font=\bfseries\footnotesize, align=center},
    linea345/.style={-, ultra thick, seni-red},
    linea138/.style={-, thick, seni-blue},
    linea230/.style={-, thick, seni-green, dashed},
    transformador/.style={-, thick, seni-gold, dashed},
    voltaje/.style={font=\footnotesize, fill=white, inner sep=1.5pt},
    genbox/.style={rectangle, draw=seni-gold, fill=seni-gold!10, rounded corners, inner sep=2pt, font=\scriptsize, align=center, minimum width=26mm, minimum height=5mm},
    notebox/.style={rectangle, draw=gray!40, fill=gray!5, rounded corners, inner sep=2pt, font=\scriptsize, align=left}
]
%% === TRIÁNGULO PRINCIPAL (anillo 138 kV) ===
\node[subestacion] (jsauri) at (0, 0.5) {JSAURI\\{\small 138 kV}};
\node[subestacion] (pala) at (-4.8, -4) {PALA\\{\small 138 kV}};
\node[subestacion] (scnu) at (4.8, -4) {SCNU\\{\small 138 kV}};
%% Anillo 138 kV (distancias reales de corredor)
\draw[linea138] (jsauri) -- node[voltaje, above left, midway] {138 kV, 70.5 km} (pala);
\draw[linea138] (jsauri) -- node[voltaje, above right, midway] {138 kV, 50.0 km} (scnu);
\draw[linea138] (pala) -- node[voltaje, below, midway] {138 kV, 22.1 km} (scnu);
%% Refuerzo 230 kV
\draw[linea230] (jsauri) -- node[voltaje, pos=0.42, right=2pt, fill=white] {230 kV} (scnu);
%% === JSAURI 345 kV (desacoplada, izquierda) ===
\node[sub345] (jsauri345) at (-8.5, 0.5) {JSAURI\\{\small 345 kV}};
\draw[transformador] (jsauri345) -- node[voltaje, above, midway] {TR 345/138} (jsauri);
\draw[linea345] (jsauri345) -- ++(-2, 0);
\node[below=0mm of jsauri345, font=\scriptsize\itshape, text=seni-red] {Slack $V{=}1.02$};
%% Generation at JSAURI 345
\node[genbox, above=8mm of jsauri345] (ptacatalina) {Punta Catalina: $2{\times}376$ MW (carbón)};
\node[genbox, above=2mm of ptacatalina] (nuclear) {Nuclear*: $4{\times}(300{-}400)$ MW};
%% === RENOVABLES (cuadrante superior derecho, fuera del anillo) ===
\node[genbox, minimum width=32mm] (renov) at (7.2, -0.8) {Renovables SCNU};
\node[genbox, below=2mm of renov] (hidro) {Hidro: 58.9 MW};
\node[genbox, below=1mm of hidro] (solar) {Solar FV: 95 MW};
\node[genbox, below=1mm of solar] (eolica) {Eólica: 34 MW};
\draw[-, thick, seni-gold, dashed] (scnu.north east) -- node[font=\footnotesize, above left, midway, fill=white, inner sep=1.5pt] {13.8 kV} (renov.south west);
%% Electrolizadores PEM (prospectivo)
\node[genbox, below=12mm of scnu, minimum width=36mm] (elec) {Electrolizadores PEM: 50 MW*};
%% === CARGA Y GENERACIÓN EN PALA ===
\node[genbox, left=2mm of pala, yshift=4mm] (carga) {Carga total: 1,645 MW};
\node[genbox, left=2mm of pala, yshift=-6mm] (fueloil) {Fuel oil: $3{\times}10.77$ MW};
%% === DISTRIBUCIÓN Y GENERACIÓN 13.8 kV ===
\node[genbox, below=12mm of pala, minimum width=38mm] (dist) {Distribución: 12.47 kV $\to$ 7.2 kV $\to$ 240/110 V};
\node[genbox, right=2mm of scnu, yshift=-6mm] (gen13) {Gen: 13.8 kV $\xrightarrow{\text{TR}}$ 138 kV};
%% === LEYENDA + NOTAS lado a lado ===
\node[notebox, below=28mm of pala, anchor=north west, minimum width=68mm] (leyenda) {
    \textbf{Leyenda:}\\[1mm]
    {\color{seni-red}\rule{6mm}{2pt}} 345 kV (transmisión)\\
    {\color{seni-blue}\rule{6mm}{1.5pt}} 138 kV (subtransmisión)\\
    {\color{seni-green}\rule{6mm}{1pt} \footnotesize---} 230 kV (refuerzo)\\
    {\color{seni-gold}\rule{6mm}{1pt} \footnotesize---} Transformador
};
\node[notebox, right=2mm of leyenda, anchor=north west, minimum width=68mm] (notas) {
    \textbf{Notas:}\\[1mm]
    *Escenario prospectivo (no operativo 2026)\\
    JSAURI 345: barra de transmisión (slack)\\
    PALA: Palamara (centro de carga metropolitano)\\
    SCNU: hub renovable + H$_2$
};
\end{tikzpicture}
}
\caption{Triángulo real de subtransmisión del anillo de 138 kV del SENI. JSAURI--PALA--SCN con longitudes de ruta calibradas de 70.5, 30.86 y 22.1 km --- las alas de corredor de 138~kV JSAURI--PALA y PALA--SCN y el enlace JSAURI--SCN de 230~kV (discontinuo), la barra de transmisión de 345 kV desacoplada a la izquierda como barra de referencia ($V = 1.02$ pu) recibiendo Punta Catalina ($2 \times 376$ MW), el hub renovable de SCN (58.9 MW hidro, 95 MW solar, 34 MW eólica) y el bloque prospectivo de electrólisis PEM de 50 MW (no operativo, 2026, con asterisco), y el centro de carga PALA (1,645 MW agregados). El catálogo $N-1$ de 215 contingencias de la batería de validación está definido sobre esta topología.}
\label{fig:topologia_seni}
\end{figure*}

% ---- Tabla III ----
\begin{table*}[!t]
\caption{Rendimiento computacional y métricas de validación del benchmark SENI-142}
\label{tab:telemetry}
\centering
\footnotesize
\begin{tabular}{@{}p{0.30\textwidth} p{0.18\textwidth} p{0.34\textwidth}@{}}
\toprule
\textbf{Métrica} & \textbf{Valor} & \textbf{Condición} \\
\midrule
Tiempo de convergencia por estado horario nacional de AC & 2.65 ms & Modelo completo de 142 barras, portátil de consumo (8 hilos) \\
Memoria residente, horizonte completo & 473 MiB & 87,360 estados horarios (520 semanas completas) \\
Residuo fasorial de acople de frontera $\delta$ & $2.788 \times 10^{-14}$ pu & Reconstrucción de Schur vs.\ estado de referencia \\
Error de regularización de límites reactivos & $5.32 \times 10^{-12}$ & Controlador sigmoidal, inyección jacobiana analítica \\
Capacidad de cribado estático (EMS) & 2.65~ms / estado & Dentro del ciclo de estimación de estado EMS de 1--5~s; cribado multiescenario de gran volumen \\
Catálogo de contingencias & 215 escenarios $N-1$ & Campaña de banco de ensayo de la Fase 2 (verificación estática post-evento) \\
Contraste espectral (IEEE 14 barras, $22\times22$) & $\kappa_2 = 117.7823$ (MARGINAL) & Control de validez externa \\
Mayor contingencia individual & Disparo de Punta Catalina, 752 MW & Factible en régimen permanente con redespacho de reserva ($P_{\mathrm{shed}} = 0$) \\
\bottomrule
\end{tabular}
\end{table*}
